\documentclass[12pt,a4paper]{letter}
\usepackage[utf8]{inputenc}
\usepackage[margin=1in]{geometry}
\usepackage{times}
\usepackage{setspace}
\onehalfspacing

% Letter details
\name{Deborah Akuoko}

\begin{document}

\begin{letter}{Professor Anna Scaife\\Professor of Radio Astronomy\\Head of Jodrell Bank Centre for Astrophysics\\University of Manchester\\Manchester\\UK\\Email: anna.scaife@manchester.ac.uk}

\date{\today}

\opening{Dear Professor Scaife,}

I hope this letter finds you well. I am writing to request your consideration in providing a letter of recommendation to support my application for the UK Global Talent visa.

I came into the UK through the DARA big data scholarship programme, of which you were then chair, through which I won a scholarship from the Newton Fund to pursue an MPhil in cancer science. During my MPhil, you helped me communicate my research by involving me in project presentation and socialisation events, especially one in Cornwall. After my MPhil, you continued to support me with a reference letter to apply to the PhD programme at the University of Edinburgh. You did not only bring me into the country, but you made sure to help me pursue my dream of a research career via PhD. Your advice and support gave me a solid foundation to thrive through the education system in the UK, especially as I came in totally new to it. Today, I find myself in need of your support once more, this time to assist with my visa application.

I have already submitted my application with a comprehensive recommendation letter from my academic supervisor, Dr. Istvan Gyongy of the University of Edinburgh, and with the full support of the University itself. However, UK Visas and Immigration has responded to my application by requesting an additional letter of recommendation from an eminent person in the UK. This official request, which I have received via email, is attached for your review along with Dr. Gyongy's recommendation letter.

The reason I am approaching you, Professor Scaife, is precisely because UKVI specifically requires a recommendation from an 'eminent' person based in the UK. By definition, an eminent person is one who has achieved distinguished recognition and standing in their field, typically evidenced by significant honours, leadership positions in prestigious institutions, and substantial contributions to their discipline. Your distinguished career exemplifies these criteria: as Professor of Radio Astronomy at the University of Manchester, Head of the Jodrell Bank Centre for Astrophysics Interferometry Centre of Excellence, co-director of Policy@Manchester, and recipient of the Jackson-Gwilt Medal by the Royal Astronomical Society in 2019 in recognition of your contributions to astrophysical instrumentation, you represent precisely the calibre of eminent person that UKVI seeks to provide such endorsements. Your standing in the UK academic and scientific community, combined with your extensive contributions to radio astronomy, big data astrophysics, and your leadership in major international projects such as the Square Kilometre Array, makes you uniquely qualified to provide the authoritative assessment that UKVI requires.

Having recently submitted my PhD thesis at the University of Edinburgh, my research has resulted in both patent applications and publication drafts that demonstrate my contributions to the field. The intellectual property and scholarly work emerging from my doctoral research reflects my competence and expertise, as well as my potential to make meaningful contributions to the academic and research landscape in the United Kingdom.


I must, with some trepidation, share that my current circumstances are marked by considerable urgency. At present, I have no source of income, and the successful outcome of this visa application is not merely a professional aspiration but a matter of personal survival. The prompt reception of the visa would go a long way to grant me financial opportunities that would enable me to help and support an ailing close relative who depends on my assistance. 

I must also acknowledge, with candour, that this matter has become more urgent for me than the completion of my PhD, for I cannot take pride in earning a prestigious degree when I find myself incapacitated as a result of having no right to work. The opportunity to continue my research and contribute to the UK's academic community through the Global Talent visa would provide not only the stability I so desperately need but also the platform to realise the potential that my supervisors and institution have recognised in my work.

I understand that requests of this nature require careful consideration, and I would be profoundly grateful for any support you might be able to offer. Your endorsement would carry significant weight in the visa application process and would represent a transformative opportunity for my career and personal circumstances.

I have enclosed the UKVI email request, the recommendation letter from Dr. Istvan Gyongy, along with documentation of my patents and publication drafts, should you wish to review the evidence of my work and the support I have already received from my academic institution.

Thank you most sincerely for your time and consideration. I would be delighted to provide any additional information you might require and would welcome the opportunity to discuss my research and aspirations with you at your convenience.

\closing{With deepest respect and gratitude,}

\signature{Deborah Akuoko}

\encl{UKVI email request for additional recommendation\\Recommendation letter from Dr. Istvan Gyongy\\Patent documentation\\Publication drafts}

\end{letter}

\end{document}

